\begin{helper}[Differentiating a mollification gives the mollification of the weak gradient]
Let $n\in\N$ and let $\Omega\subseteq\R^n$ be measurable, and suppose that a
function vanishing off $A\subseteq B$ has the same integral over $B$ as over
$A$.  Throughout, $\operatorname{supp}\rho$ denotes the \emph{closed} support
of $\rho$, $\partial_iv(x)=dv(x)[e_i]$, $1_Kf=K$-indicator of $f$, and
$u*\rho$ is the convolution against Lebesgue measure with scalar multiplication
as the pairing.  Let $K_3\subseteq\R^n$, let $f\colon\R^n\to\R$, let
$G_f\colon\R^n\to\R^n$, let $\rho\colon\R^n\to\R$, let $d\in\R$ and let
$x\in\R^n$.  Assume
\begin{enumerate}
\item[(H1)] $f$ has weak gradient $G_f$ on $\Omega$
  (\noderef{HasWeakGradientOn});
\item[(H2)] $1_{K_3}f$ is locally integrable;
\item[(H3)] $\rho$ is smooth with compact support and
  $\operatorname{supp}\rho=\overline{B_d(0)}$;
\item[(H4)] $\overline{B_d(x)}\subseteq K_3$ and
  $\overline{B_d(x)}\subseteq\Omega$.
\end{enumerate}
Then for every coordinate $i$,
\[
\partial_i\bigl((1_{K_3}f)*\rho\bigr)(x)
\ =\ \bigl((1_{K_3}G_f^{\,i})*\rho\bigr)(x).
\]
Nothing here refers to a Sobolev exponent, to a composition or to an outer
function: it is a statement about one weakly differentiable function and one
mollifier.
\end{helper}
\begin{proof}
\emph{Step 1: differentiation under the convolution.}  By (H3), $\rho$ is
smooth with compact support, so $\partial_i\rho$ is continuous with compact
support; together with the local integrability of $1_{K_3}f$ from (H2) this
makes both convolutions $(1_{K_3}f)*\rho$ and $(1_{K_3}f)*\partial_i\rho$ exist,
and differentiation passes under the integral sign (Preamble), giving
\[
\partial_i\bigl((1_{K_3}f)*\rho\bigr)(x)
=\bigl((1_{K_3}f)*\partial_i\rho\bigr)(x)
=\int_{\R^n}(1_{K_3}f)(y)\,\partial_i\rho(x-y)\,dy .
\]

\emph{Step 2: the shifted mollifier is an admissible test function.}  Fix $x$
as in (H4) and put $\varphi(y)=\rho(x-y)$.  Then $\varphi$ is smooth, and
$\operatorname{supp}\varphi=x-\operatorname{supp}\rho=\overline{B_d(x)}$ by
(H3), which is contained in $\Omega$ by (H4); so $\varphi$ is a smooth
compactly supported test function on $\Omega$.  Moreover
$\partial_{y_i}\varphi(y)=-\partial_i\rho(x-y)$ by the chain rule.

\emph{Step 3: the weak-gradient identity.}  Applying (H1) to $\varphi$,
\[
\int_{\Omega}f(y)\,\partial_{y_i}\varphi(y)\,dy
\ =\ -\int_{\Omega}G_f^{\,i}(y)\,\varphi(y)\,dy ,
\]
that is, $-\int f(y)\partial_i\rho(x-y)dy=-\int G_f^{\,i}(y)\rho(x-y)dy$, and
the two minus signs cancel:
$\int f(y)\partial_i\rho(x-y)dy=\int G_f^{\,i}(y)\rho(x-y)dy$.

\emph{Step 4: inserting and removing the indicator.}  Every integrand above
vanishes off $\overline{B_d(x)}$, because $\varphi$ and each
$\partial_{y_i}\varphi$ do, and $\overline{B_d(x)}\subseteq K_3$ by (H4).  The
hypothesis on integrals of functions vanishing off a subset therefore lets the
domain be cut down to $\overline{B_d(x)}$ and the factor $1_{K_3}$ inserted or
deleted freely on both sides; in particular
$\int(1_{K_3}f)(y)\partial_i\rho(x-y)dy=\int f(y)\partial_i\rho(x-y)dy$ and
$\int G_f^{\,i}(y)\rho(x-y)dy=\int(1_{K_3}G_f^{\,i})(y)\rho(x-y)dy$.

\emph{Conclusion.}  Chaining Steps 1, 3 and 4,
\[
\partial_i\bigl((1_{K_3}f)*\rho\bigr)(x)
=\int(1_{K_3}G_f^{\,i})(y)\,\rho(x-y)\,dy
=\bigl((1_{K_3}G_f^{\,i})*\rho\bigr)(x),
\]
which is the assertion.
\end{proof}
